\begin{abstract}
	均质压燃（HCCI）燃烧作为一种能实现高效低污染燃烧方式，能够使发动机同时具备较高的经济性、动力性和排放性能，可有效降低发动机NOx和碳烟排放。然而HCCI燃烧的一个显著特点是燃烧着火时刻和燃烧放热率主要受化学动力学控制，因此，发动机结构参数和运行工况的改变将显著影响HCCI燃烧着火和燃烧放热过程。

	本文以二甲醚（DME）为燃料，对HCCI燃烧着火和燃烧放热过程进行了数值研究。研究中采用了美国Lawrence Livermore国家实验室提出的DME详细化学动力学反应机理和该实验室开发的HCT化学动力学计算程序，对DME均质压燃燃烧过程进行了详细模拟。新的DME动力学机理包含79个物种和399个反应。为考虑壁面传热的影响，在HCT程序中加入了壁面传热模型。用该方法研究了压缩比、燃空当量比、进气加热、发动机转速、EGR和燃料添加剂等对HCCI着火和燃烧的影响。

	研究结果表明，DME的HCCI燃烧过程由低温反应放热阶段和高温反应放热阶段组成。还发现增加压缩比、当量比、进气温度和H2O2、H2或CO的含量使着火提前。增加发动机转速、采用冷EGR和CH4或CH3OH的含量将推迟着火。

	\keywords{均质压燃，化学动力学，数值模拟，二甲醚，EGR}
\end{abstract}

\begin{enabstract}
	HCCI (Homogeneous Charge Compression Ignition) combustion has advantages in terms of efficiency and reduced emission. HCCI combustion can not only ensure both the high economic and dynamic quality of the engine, but also efficiently reduce the NOx and smoke emission. Moreover, one of the remarkable characteristics of HCCI combustion is that the ignition and combustion process are controlled by the chemical kinetics, so the HCCI ignition time can vary significantly with the changes of engine configuration parameters and operating conditions.

	In this work numerical scheme for the ignition and combustion process of DME homogeneous charge compression ignition is studied. The detailed reaction mechanism of DME proposed by American Lawrence Livermore National Laboratory (LLNL) and the HCT chemical kinetics code developed by LLNL are used to investigate the ignition and combustion processes of an HCCI engine fueled with DME. The new kinetic mechanism for DME consists of 79 species and 399 reactions. To consider the effect of wall heat transfer, a wall heat transfer model is added into the HCT code. By this method, the effects of the compression ratio, the fuel-air equivalence ratio, the intake charge heating, the engine speed, EGR and fuel additive on the HCCI ignition and combustion are studied.

	The results show that the HCCI combustion fueled with DME consists of a low temperature reaction heat release period and a high temperature reaction heat release period. It is also founded that increasing the compression ratio, the equivalence ratio, the intake charge temperature and the content of H2O2, H2 or CO cause advanced ignition timing. Increasing the engine speed, adoption of cold EGR and the content of CH4 or CH3OH will delay the ignition timing.

	\enkeywords{HCCI, chemical kinetics, numerical simulation, DME, EGR}
\end{enabstract}
